\section{Texto e listas}

\begin{frame}[fragile]{Formatação de texto}
\begin{columns}[T]
\begin{column}{0.52\textwidth}
\begin{lstlisting}
Texto em \textbf{negrito}.
Texto em \textit{italico}.
Uma palavra \emph{enfatizada}.
O sistema \LaTeX{}.
Uma quebra de linha\\aqui.
\end{lstlisting}
\end{column}
\begin{column}{0.44\textwidth}
Texto em \textbf{negrito}.\\
Texto em \textit{itálico}.\\
Uma palavra \emph{enfatizada}.\\
O sistema \LaTeX{}.\\
Uma quebra de linha\\aqui.
\end{column}
\end{columns}

Use chaves para indicar exatamente qual trecho recebe a formatação.
\end{frame}

\begin{frame}[fragile]{Caracteres reservados e comentários}
\begin{lstlisting}
% Esta linha nao aparece no PDF.
Taxa de acerto: 95\%
Arquivo: dados\_treino.csv
Comandos com \{chaves\}
Parametro \#1 custa R\$ 10.
\end{lstlisting}

\begin{itemize}
    \item \verb|%| inicia um comentário.
    \item \verb|\%|, \verb|\_|, \verb|\#|, \verb|\$|, \verb|\{| e \verb|\}| imprimem caracteres reservados.
    \item Uma linha em branco inicia outro parágrafo.
\end{itemize}
\end{frame}

\begin{frame}[fragile]{Listas com marcadores}
\begin{columns}[T]
\begin{column}{0.52\textwidth}
\begin{lstlisting}
\begin{itemize}
  \item Grande quantidade de imagens
  \item Imagens em diferentes estados
  \begin{itemize}
    \item brilho
    \item contraste
  \end{itemize}
\end{itemize}
\end{lstlisting}
\end{column}
\begin{column}{0.44\textwidth}
\begin{itemize}
  \item Grande quantidade de imagens
  \item Imagens em diferentes estados
  \begin{itemize}
    \item brilho
    \item contraste
  \end{itemize}
\end{itemize}
\end{column}
\end{columns}
\end{frame}

\begin{frame}[fragile]{Listas numeradas}
\begin{columns}[T]
\begin{column}{0.52\textwidth}
\begin{lstlisting}
\begin{enumerate}
  \item Preparar os arquivos
  \item Compilar o documento
  \item Revisar o PDF
\end{enumerate}
\end{lstlisting}
\end{column}
\begin{column}{0.44\textwidth}
\begin{enumerate}
  \item Preparar os arquivos
  \item Compilar o documento
  \item Revisar o PDF
\end{enumerate}
\end{column}
\end{columns}

Troque \verb|itemize| por \verb|enumerate| quando a ordem dos itens importar.
\end{frame}
